\documentclass[11pt,a4paper]{article}
\usepackage[T1]{fontenc}
\usepackage{lmodern,amsmath,amssymb,amsthm,microtype}
\usepackage[margin=26mm]{geometry}
\usepackage[colorlinks=true,linkcolor=blue,urlcolor=blue]{hyperref}
\newtheorem{theorem}{Theorem}[section]
\newtheorem{lemma}[theorem]{Lemma}
\newtheorem{corollary}[theorem]{Corollary}
\newcommand{\Q}{\mathbb Q}\newcommand{\Z}{\mathbb Z}
\title{Elliptic index divisibility and amplified congruences}
\author{}\date{17 November 2026}
\begin{document}\maketitle

\begin{abstract}
For a fixed rational congruent-number curve, a reciprocal square
abscissa supplies a balanced pole test at the odd denominator primes
of an elliptic multiple. Under weak rational $abc$ with any fixed
exponent strictly between one and two, an index progression ensures
that every solution outside a fixed finite exceptional set forces
divisibility of the two elliptic indices. Conversely, every nonzero
integer ratio occurs for infinitely many index pairs satisfying the
complete intrinsic-pole preparation. Adding a local logarithmic-ratio
test preserves completeness. Formal-logarithm truncation amplifies
the resulting congruences to any fixed multiplicity without a new free
coordinate. A separate height condition is required to identify an
external rational input with the index ratio.
\end{abstract}

\section{The source}\label{sec:source}

Fix $d\in\Z_{>0}$ and
\[
                         E_d:y^2=x^3-d^2x.
\]
Let $P\in4E_d(\Q)$ have infinite order. Choose its sign and replace
it by a fixed positive multiple so that it belongs to the formal group at two with
parameter $\zeta(P)=-x(P)/y(P)$ of valuation at least four.
Write
\[
              H=[h]P,\qquad U=[j]P,\qquad h,j\in\Z\setminus\{0\}.
\]
These index representations are hypotheses of the arithmetic statement.
For a rank-one curve they may be enforced by restricting both points
to a fixed multiple of $E_d(\Q)$. This restriction is a finite
isogeny cover and introduces no further free point coordinate.

Let $q_H\in\Q^\times$ satisfy $q_H^2=x(H)$. There is a positive
choice $q_0=\sqrt{x(P)}\in\Q$. Choose a fixed positive rational
$\lambda_0$ with
\[
                  v_2(\lambda_0)=2,\qquad v_7(\lambda_0)=-1,
\]
and sufficiently large real absolute value, as specified below. Set
$c=q_0\lambda_0^2$ and define
\begin{equation}\label{eq:weight}
                       u=\frac c{q_H},\qquad w=\frac{u}{(1+u)^2}.
\end{equation}
Use the test $C_1=f(U)=1/x(U)$. Optionally, for an external rational
input $X$, also use
\[
                 C_2=\frac{\zeta(U)}{\zeta(H)}-X.
\]
For either test list put
\begin{equation}\label{eq:detector}
 D_i=\frac{C_i^2}{1+uC_i^2},\qquad
 W=\prod_i\left(1+\frac{D_i^2}{u}\right),\qquad
 b=4wW^2,\qquad \eta=7b.
\end{equation}
Retain all nonzero denominators. Impose the intrinsic equation
\begin{equation}\label{eq:pole}
 Y^2=\eta(\eta+v^2)^2s^4
             -\eta(\eta-v^2)^2(1-\eta)^2,
                \qquad v,s,Y\in\Q.
\end{equation}
The intrinsic pole theorem \cite{Intrinsic} implies from this equation
that $\eta$ has square reduced denominator.

\section{Soundness}

Let $S$ contain $2,7$, the bad reduction primes, and the numerator
and denominator primes of $c$. Write
\[
       x([n]P)=A_n^2/B_n^2,
       \qquad A_n,B_n>0,\quad\gcd(A_n,B_n)=1.
\]
Fix $1<\kappa<2$, choose an integer
$R_\kappa>1/(2/\kappa-1)$, and put
\[
                 N_\kappa=\prod_{\substack{p\le R_\kappa\\p\text{ prime}}}p.
\]
Restrict $h\equiv1\pmod{N_\kappa}$. Under weak rational $abc$
with exponent $\kappa$, the coprime-index version of the simple
primitive-prime theorem of \cite{Simple} gives $n_0$ such that
$B_{|h|}$ has a primitive prime of valuation one outside $S$
whenever $|h|\ge n_0$. Its proof uses the strict inequality
\[
        \sum_{\substack{\ell\mid h\\\ell\text{ prime}}}\ell^{-2}
                         \le1/R_\kappa<2/\kappa-1.
\]

On a rank-one curve, suppose $P=L G$ for a generator $G$
modulo torsion, where $L$ kills torsion and imposes the retained
subgroup. The required coset is enforced by
\[
                 H-P\in[LN_\kappa]E_d(\Q),
                 \qquad U\in[L]E_d(\Q).
\]
The factor $L$ in the first condition is essential. Both conditions
are finite isogeny covers over the retained image points.

\begin{theorem}\label{thm:sound}
Assume weak rational $abc$ with a fixed exponent $1<\kappa<2$,
$h\equiv1\pmod{N_\kappa}$ and $|h|\ge n_0$.
Every rational solution of \eqref{eq:pole}
in the guarded source satisfies $h\mid j$. The assertion remains
valid when any additional rational tests are included in
\eqref{eq:detector}.
\end{theorem}
\begin{proof}
Choose a simple primitive prime $p\notin S$ of $B_{|h|}$.
Since $c$ is a unit there, $v_p(u)=1$.
More generally, when $\nu=v_p(u)$ is odd,
\[
                  v_p(D_i)=\max(2v_p(C_i),-\nu).
\]
The two terms of $1+uC_i^2$ have distinct valuations, which proves
this formula, with the convention $v_p(0)=+\infty$.
Every factor $1+D_i^2/u$ has valuation at most zero: its two terms
have distinct valuations because $2v_p(D_i)-\nu$ is odd.

If $U$ is not formal at $p$, then $v_p(C_1)\le0$, so
$v_p(D_1)\le0$ and $v_p(W)\le-\nu=-1$. Thus
\[
                  v_p(b)=\nu+2v_p(W)<0
\]
is odd. The same is true of $7b$ because $p\ne7$, contradicting
\eqref{eq:pole}. Therefore $U$ is formal at $p$.
The order of $\widetilde P$ modulo $p$ is exactly $|h|$ by
primitivity. Since $[j]\widetilde P=O$, it follows that $h\mid j$.
\end{proof}

The finite restriction $|h|\ge n_0$ requires no free coordinate.
After the cyclic subgroup has been imposed, it can be expressed by
the nonzero guard
\[
                 \prod_{1\le r<n_0}\bigl(x(H)-x([r]P)\bigr)\ne0.
\]
The affine point guards already exclude $h=0$.

\section{Completeness for every integer ratio}

\begin{theorem}\label{thm:complete}
The fixed constant $c$ may be chosen so that the following holds.
For every nonzero integer $m$, there are infinitely many nonzero
integers $h\equiv1\pmod{N_\kappa}$ with $|h|\ge n_0$ for which,
setting $j=mh$ and
choosing $q_H>0$, the input $b$ in \eqref{eq:detector} satisfies
\[
 0<b<1/7,\qquad b\in\Q_2^{\times2},\qquad v_2(b)\ge4,
 \qquad v_7(b)\ge0,
\]
and has square reduced denominator. This also holds with the
optional test $C_2$ when $X=m$. Consequently the intrinsic
theorem supplies rational solutions of \eqref{eq:pole}.
\end{theorem}
\begin{proof}
Choose $\lambda_0$ as above with $\lambda_0^2>100$.
Take $h\equiv1\pmod{N_\kappa}$ in a sufficiently
small fixed congruence neighborhood of one so that $H=[h]P$
is close to $P$ at every prime in $S$. Choose the rational
square-root branch $q_H$ close to $q_0$ at these primes. This
branch is coherent: since $P=2R$ for some fixed rational point
$R$, the rational function
\[
        q([h]R)=\frac{x([h]R)^2+d^2}{2y([h]R)}
\]
is a square root of $x([h]P)$; changing the initial sign of $R$
if necessary makes its value $q_0>0$ at $h=1$. Congruences
may also be imposed on $[h]R$.
At the odd primes of $S$, closeness ensures
$v_p(u)=v_p(\lambda_0^2)$, which is even. At two and seven, refine it
so that $q_0/q_H$ is a local square. Thus
\[
 u\in\Q_2^{\times2},\quad v_2(u)=4,\qquad
 u\in\Q_7^{\times2},\quad v_7(u)=-2.
\]

The formal group at two gives $v_2(\zeta(H))\ge2$ and
$v_2(\zeta(U))\ge v_2(\zeta(H))$, since $U=[m]H$ and
$m\in\Z$. Hence $v_2(C_1)\ge4$. Also the integral formal
multiplication series gives
\[
       \frac{\zeta([m]H)}{\zeta(H)}-m
                      \in\zeta(H)\Z_2,
\]
so $v_2(C_2)\ge2$ when this test is present. Therefore all
factors of $W$ are two-adic units congruent to one modulo two.
It follows that $b$ is a two-adic square and $v_2(b)=6$.

At seven the terms of $1+uC_i^2$ cannot cancel: in the
equal-valuation case $uC_i^2$ has square residue, whereas
$-1$ is not a square modulo seven. Thus the displayed valuation
formula for $D_i$ remains valid even though $v_7(u)=-2$.
It gives $v_7(D_i)\ge2$, so $W\equiv1\pmod7$ and $v_7(b)=2$.

Consider any other prime. If $v_p(u)$ is even, then
$v_p(b)=v_p(4)+v_p(u)-2v_p(1+u)+2v_p(W)$ is even.
If $\nu=v_p(u)<0$ is odd, the balanced valuation formula gives
$v_p(D_i^2/u)\ge-3\nu>0$, hence $W\equiv1\pmod p$ and
$v_p(b)=-\nu>0$.
It remains to consider odd positive $\nu$ outside $S$.
Here $H$ is formal and $v_p(\zeta(H))=\nu$, while $U=[m]H$.
Thus $v_p(C_1)\ge2\nu$ and, if present,
$v_p(C_2)\ge\nu$ by the same integral formal multiplication
series. Every factor of $W$ is a unit congruent to one, and
$v_p(b)=\nu>0$. This proves square denominator at all primes.

Finally take $[h]R$ in a sufficiently small real neighborhood
of $R$. Then $q_H>0$ and $u$ remains larger than $50$.
All tests are real, and
\[
       0\le D_i<1/u,\qquad
       1\le W\le(1+u^{-3})^2.
\]
Consequently $0<b\le4u^{-1}(1+u^{-3})^4<1/7$.
Every fixed arithmetic progression of multiples of a nontorsion
real elliptic point is dense in the appropriate connected real
component. Therefore this real neighborhood is reached for
infinitely many arbitrarily large $h$ in the specified congruence
class. The nonzero affine guards exclude at most finitely many
of these indices. Apply the intrinsic completeness theorem to
$\eta=7b$.
\end{proof}

\section{A monic algebra for the finite isogeny covers}

Let $R_n(Z)=\mathcal A_n(Z)/\mathcal B_n(Z)$ be the rational
function giving the abscissa of multiplication by a positive integer
$n$ on $E_d$. Choose coprime numerator and denominator with
$\mathcal A_n$ monic. The degrees are $n^2$ and $n^2-1$,
respectively; the leading coefficient of $\mathcal B_n$ is $n^2$.
These follow either from the division polynomials or from the degree
of $[n]$ and its formal expansion at infinity; see
\cite[Chapter III]{AEC}.

\begin{lemma}\label{lem:isogeny-algebra}
Let $K$ be a characteristic-zero field and let
$T_0=(x_0,y_0)\in E_d(K)$ with $y_0\ne0$. The monic algebra
\[
 \mathcal C=K[Z]/\bigl(\mathcal A_n(Z)-x_0\mathcal B_n(Z)\bigr)
\]
is finite etale of rank $n^2$. Both $\mathcal B_n(Z)$ and
$R_n'(Z)$ are units in $\mathcal C$. The formulas
\[
                  x_R=Z,\qquad y_R=\frac{ny_0}{R_n'(Z)}
\]
give a point $R\in E_d(\mathcal C)$ satisfying $[n]R=T_0$.
Every $K$-rational preimage of $T_0$ is obtained uniquely from
a $K$-algebra character of $\mathcal C$.
\end{lemma}
\begin{proof}
Coprimality of $\mathcal A_n,\mathcal B_n$ makes
$\mathcal B_n$ a unit in the quotient. Pullback of the invariant
differential under $[n]$ gives
\[
 y([n]R)=\frac{y(R)R_n'(x(R))}{n},\qquad
 (R_n(Z))^3-d^2R_n(Z)
     =\frac{(Z^3-d^2Z)(R_n'(Z))^2}{n^2}.
\]
In $\mathcal C$ the left side of the second identity is
$y_0^2\ne0$. It follows that $R_n'(Z)$ is a unit, and the
displayed formulas for $x_R,y_R$ satisfy the curve equation
and map to $T_0$. Moreover, modulo the defining polynomial,
\[
 (\mathcal A_n-x_0\mathcal B_n)'
                         =\mathcal B_n R_n'
\]
is a unit, so the algebra is etale. Its rank is the degree
$n^2$ of the monic polynomial. The ordinate is uniquely
determined by the abscissa and the chosen target ordinate,
which proves the assertion about characters.
\end{proof}

Apply the lemma with $n=LN_\kappa$ and target $H-P$, and
with $n=L$ and target $U$. Retain the affine target guards and
their nonzero ordinates. On the intended source, the guard
$x(H)\ne x(P)$ excludes $h=\pm1$ and ensures that both
sign characters of $H$ give a nontorsion target $H-P$.
The finite exceptional-index guard may always include this
factor. The isogeny part has rank $(LN_\kappa)^2L^2$ over the
image-point algebra and introduces no new free coordinate.
The same construction over a finite etale base algebra applies
componentwise when the target guards are units.

\section{Finite coordinates and guards}\label{sec:scope}

On a rational elliptic curve each point uses one retained rational
abscissa coordinate and finite algebraic roots. The pole contributes
two free parameters, while its ordinate is a finite quadratic root.
Finite isogeny covers and the displayed nonzero guards introduce no
additional free parameters in the finite-algebra encoder.
Lemma~\ref{lem:isogeny-algebra} implements the required isogeny
cosets by monic finite algebras.

For the square-abscissa chart, retain $q_H,q_U$ freely and use
the monic roots $l_H^2=q_H^4-d^2$, $l_U^2=q_U^4-d^2$, with
$y(H)=q_Hl_H$ and $y(U)=q_Ul_U$. At the prepared rational
values all these point characters are real, while $u=c/q_H>50$
is independent of their signs. Thus every balanced denominator,
$W$, and $b$ is nonzero at every point character. A finite isogeny
cover does not change the images of these point characters.
The intrinsic density theorem allows the pole parameters to avoid
the radicand and division guards simultaneously at every character.

The theorem establishes an integer \emph{index ratio}. With the
optional comparison test, it supplies local congruences between
$X$ and $j/h$. Identifying these two rational numbers requires
an additional argument. No defining-variable bound for
$\Z\subset\Q$ is asserted here.

\section{The amplified detector}

Use the notation, fixed weak exponent $1<\kappa<2$, and cyclic
point coset of Sections~\ref{sec:source}--\ref{sec:scope} on
$E_d:y^2=x^3-d^2x$, with $P\in4E_d(\Q)$,
$H=[h]P$, $U=[j]P$, and $q_H^2=x(H)$.
Fix an integer $T\ge1$ and put $k=8T-7$.
Let $\Lambda_T(Z)\in\Q[Z]$ be the degree-at-most-$T$ truncation
of the formal logarithm in the parameter $\zeta=-x/y$, normalized
to begin with $Z$. In particular $\Lambda_1(Z)=Z$.
For this curve the coefficients are explicitly
\[
 \log_E(Z)=\sum_{r\ge0}
      \frac{(-1)^r\binom{2r}{r}d^{2r}}{4r+1}Z^{4r+1}.
\]
Indeed, with $w=-1/y$ the curve equation is
$w=Z^3-d^2Zw^2$, and its invariant differential is
$dZ/\sqrt{1+4d^2Z^4}$.
Replace the weight and test list by
\[
 u=\left(\frac c{q_H}\right)^k,\qquad
 C_1=f(U)^T,\qquad
 C_2=\frac{\Lambda_T(\zeta(U))}{\Lambda_T(\zeta(H))}-X.
\]
Keep the balanced definitions
\[
 D_i=\frac{C_i^2}{1+uC_i^2},\quad
 W=\prod_{i=1}^2\left(1+\frac{D_i^2}{u}\right),\quad
 b=\frac{4uW^2}{(1+u)^2},\quad\eta=7b,
\]
and the intrinsic pole equation
\[
 Y^2=\eta(\eta+v^2)^2s^4
              -\eta(\eta-v^2)^2(1-\eta)^2.
\]
Retain the previous guards and also
$\Lambda_T(\zeta(H))\ne0$. The finite exceptional set $S$
contains $2,7$, the bad reduction and coefficient primes, and
every prime at most $T+1$.

\begin{lemma}\label{lem:truncation}
At every good prime $p\notin S$, if $H$ is formal with
$v_p(\zeta(H))=e>0$ and $U=[m]H$ for $m\in\Z\setminus\{0\}$,
then
\[
 v_p\left(
 \frac{\Lambda_T(\zeta(U))}{\Lambda_T(\zeta(H))}-m
 \right)\ge Te.
\]
\end{lemma}
\begin{proof}
The invariant differential has integral formal coefficients, so
the coefficient of $Z^r$ in its logarithm has valuation at least
$-v_p(r)$. For $r\ge T+1$ and $p>T+1$ one has
$r-1-v_p(r)\ge T$: this is immediate when $p\nmid r$, while
$p^a-1-a\ge T$ when $a=v_p(r)\ge1$.
Consequently
\[
 \frac{\Lambda_T(z)}{\log_E(z)}\in1+p^{Te}\Z_p
                         \qquad(v_p(z)=e).
\]
Both the logarithm and its truncation have valuation $e$.
Now use $\log_E(U)=m\log_E(H)$ and
$v_p(\zeta(U))\ge e$. Multiplication by the integral number
$m$ does not decrease the stated error valuation.
\end{proof}

\begin{theorem}\label{thm:amplified}
Assume weak rational $abc$ with the fixed exponent $1<\kappa<2$,
retain $h\equiv1\pmod{N_\kappa}$, and exclude the
same finite exceptional set of indices as in Sections~\ref{sec:source}--\ref{sec:scope}, enlarged
for the present $S$. Every rational pole solution satisfies
$h\mid j$. Set $m=j/h$.
If $p\notin S$ satisfies $v_p(B_{|h|})=1$, then
\[
                              v_p(X-m)\ge T.
\]
For every nonzero integer $m$, there are infinitely many $h$
in the prescribed index progression
with $j=mh$ and $X=m$ for which the complete intrinsic-pole
preparation holds. The initial choice $v_2(\zeta(P))\ge4$
suffices simultaneously for every $T$. The amplification introduces no new
free parameter.
\end{theorem}
\begin{proof}
At a simple primitive denominator prime of $H$, $v_p(u)=k$.
If $U$ is not formal, then $v_p(C_1)\le0$. The same balanced
valuation argument as in Sections~\ref{sec:source}--\ref{sec:scope} gives a negative odd
valuation of $b$, which is impossible. Thus $h\mid j$.

At any simple denominator prime, not necessarily primitive,
$v_p(u)=k$. Every factor of $W$ has valuation at most zero.
If $v_p(D_2)\le k/4$, the second factor has negative valuation
and
\[
 v_p(b)=k+2v_p(W)\le4v_p(D_2)-k<0;
\]
the last strict inequality follows because the left side of
the preceding bound is odd. Hence $v_p(D_2)>k/4$.
Its positivity implies $v_p(D_2)=2v_p(C_2)$, so
\[
                      v_p(C_2)>k/8=T-7/8,
\]
and the integral valuation is at least $T$.
Lemma~\ref{lem:truncation} now gives $v_p(X-m)\ge T$.

For completeness take $j=mh$ and $X=m$. At any positive odd
denominator depth $e=v_p(q_H^{-1})$ outside $S$, one has
$v_p(u)=ke$, $v_p(C_1)\ge2Te$ and $v_p(C_2)\ge Te$.
The first factor of $W$ is a unit. The second has valuation
at least $\min(0,(4T-k)e)$, and therefore
\[
 v_p(b)\ge\min(k,8T-k)e=\min(8T-7,7)e>0.
\]
At negative odd valuations of $u$, the balanced transformation
makes $W$ a unit and $b$ has positive valuation. Even
valuations of $u$ always give even valuations of $b$.

At the fixed places choose $c=q_0\lambda_0^2$ and the progression of $h$
as in Sections~\ref{sec:source}--\ref{sec:scope}. Then $u$ is a square at two and seven,
with valuations $4k$ and $-2k$, respectively. All coefficients
of the displayed formal logarithm are two-adic integers,
because the denominators $4r+1$ are odd. Therefore the truncation
estimate $v_2(C_2)\ge T v_2(\zeta(H))$ holds at two as well,
uniformly in $T$ and the integer $m$.
Choose the initial $P$ with $v_2(\zeta(P))\ge4$.
Then $v_2(C_1)\ge8T$ and $v_2(C_2)\ge4T$, so both
balanced denominators are units. The first factor of $W$ is
a unit since $4v_2(C_1)-4k\ge28>0$. For the second,
possible cancellation can only increase its valuation. Hence
\[
 v_2(b)\ge2+4k+2\min(0,16T-4k)
             =2+\min(4k,28)\ge6.
\]
The input $b$ is a two-adic square because $u$ is.
At seven, the non-cancellation argument using the nonsquareness
of $-1$ gives $W\equiv1\pmod7$ and $v_7(b)=2k$.
The other fixed places have even valuation of $u$.

Real preparation is unchanged: when $c/q_H>50$, also $u>50$,
and $0\le D_i<1/u$ uniformly in both tests. Hence $0<b<1/7$.
The logarithm guard excludes finitely many points, and the
usual real density in a fixed arithmetic progression supplies
infinitely many admissible $h$. The intrinsic theorem gives
the rational pole witnesses and simultaneous norm-guard avoidance.
\end{proof}

\section{The exact remaining height condition}

Write $H_H=h(q_H)$ and let $L_{h,S}$ be the product of the
simple denominator primes of $H$ outside $S$. The refined
primitive-denominator argument of \cite{Simple} gives
\[
 \log L_{h,S}\ge\left(\frac2\kappa-1\right)H_H-O(1),
\]
under the fixed weak exponent $1<\kappa<2$.
If $X\ne m$, Theorem~\ref{thm:amplified} implies
\begin{equation}\label{eq:remaining}
 T\left(\frac2\kappa-1\right)H_H
           \le h(X)+\log(1+|m|)+O_T(1).
\end{equation}
Indeed $L_{h,S}^T$ divides the reduced numerator of $X-m$,
whose absolute value is at most
$\exp(h(X))(1+|m|)$.

Consequently any additional source condition giving
\[
       h(X)+\log(1+|m|)\le C H_H+O(1),
       \qquad C<T(2/\kappa-1),
\]
would bound $H_H$ for every false input. Excluding the
corresponding finite set of elliptic indices would force
$X=m\in\Z$. For a true input $X=m$, the numerical inequality
itself is compatible with arbitrarily large $h$.
What is missing is a complete, low-coordinate existential
encoding of such a growth condition. A real inequality alone
does not bound the rational height of $X$, and a finite
collection of congruences does not bound the integer $m$.

\begin{thebibliography}{9}
\bibitem{AEC}J.~H.~Silverman, \emph{The Arithmetic of Elliptic Curves},
second edition, Graduate Texts in Mathematics 106, Springer, 2009,
Chapter III.
\bibitem{Simple}\emph{Simple primitive denominator primes under weak
$abc$}, companion paper, \nolinkurl{simple_primitive_denominators_20261117_v60.tex}.
\bibitem{Intrinsic}\emph{A two-coordinate pole test of coefficient
height four}, companion paper, \nolinkurl{intrinsic_pole_20261117_v56.tex}.
\end{thebibliography}
\end{document}
